\chapter{The Hawkes process and its cluster representation}
\label{ch:hawkes}

\begin{motivation}
The Hawkes process is the model we actually believe in: events beget events.
Understanding it as a \emph{branching} process --- a forest of family trees --- is
what makes its mean computable (Ch.~\ref{ch:mbpp}), its simulation easy (this
chapter), and its failure to be Poisson obvious.
\end{motivation}

\section{Definition}

\begin{model}[Hawkes process]\label{mod:ch3-hawkes}
The intensity is
$\lambda(t)=s(t)+\int_0^{t^-}\phi(t-u)\dd N(u)=s(t)+\sum_{t_i<t}\phi(t-t_i)$,
with exogenous rate $s\ge0$ and triggering kernel $\phi\ge0$. The \emph{exponential
kernel} $\phi(t)=\kappa\theta e^{-\theta t}$ has $\int_0^\infty\phi=\kappa$ (the
\emph{branching ratio}) and decay rate $\theta$.
\end{model}

\begin{whybox}[: the exponential kernel]
We default to $\phi=\kappa\theta e^{-\theta t}$ because it is the one kernel whose
self-convolutions telescope (Ch.~\ref{ch:solving}), turning every later equation
into a closed form or a tiny ODE. Heavy-tailed (power-law) kernels are more
realistic for social media but cost us the closed form --- a trade we revisit in
Ch.~\ref{ch:solvers}.
\end{whybox}

\section{The cluster (branching) representation}

\begin{theorem}[Hawkes--Oakes]\label{thm:ch3-cluster}
The Hawkes process equals the following branching construction: immigrants arrive
as a Poisson process of intensity $s$; each event at $u$ independently spawns
offspring as a Poisson process of intensity $\phi(\cdot-u)$ on $(u,\infty)$; the
process is the superposition of all immigrants and descendants.
\end{theorem}

\begin{intuitionbox}
Self-excitation, made literal: every event founds a family, and the process is a
\emph{forest of independent family trees} rooted at the immigrants. The mean number
of direct children per event is $\int\phi=\kappa$, so a family is a Galton--Watson
tree with mean offspring $\kappa$: finite on average iff $\kappa<1$
(\emph{subcritical}), with expected total size $1/(1-\kappa)$.
\end{intuitionbox}

This picture is also the simulator: sample immigrants from $s$, then recursively
sample each event's offspring. The package does exactly this.

\begin{lstlisting}[caption={Cluster simulation (from \code{ic\_simulate.py}).}]
def simulate_separable_hawkes(exogenous, kappa, theta, T, seed=None):
    rng = np.random.default_rng(seed)
    immigrants = sample_immigrants(exogenous, T, rng)   # from s(t)
    offspring = []
    current = list(immigrants)
    while current:                       # breadth-first over generations
        nxt = []
        for p in current:
            n = rng.poisson(kappa)       # mean kappa children per event
            child = p + rng.exponential(1.0/theta, size=n)   # Exp(theta) delays
            child = child[child <= T]
            offspring.extend(child); nxt.extend(child)
        current = nxt
    return np.sort(immigrants), np.sort(np.array(offspring))
\end{lstlisting}

\begin{whybox}[: keep immigrants and offspring separate]
Returning the two streams separately lets us later reconstruct \emph{any} of the
six observation regimes (event-time vs.\ interval-censored, exogenous vs.\
endogenous) from one simulation --- essential for the experiments in
Ch.~\ref{ch:fitting}.
\end{whybox}

\section{Why it is not Poisson, and the mean-intensity equation}

Because $\lambda$ depends on the realised history, the compensator
$\Lambda=\int\lambda$ is \emph{random}; by Watanabe (Thm~\ref{thm:ch2-watanabe})
the Hawkes process is therefore not Poisson and has dependent increments. But its
\emph{mean} is deterministic, and satisfies a clean equation.

\begin{proposition}[mean-intensity equation]\label{prop:ch3-mean}
$\xi(t)=\E[\lambda(t)]$ satisfies
$\xi(t)=s(t)+\int_0^t\phi(t-u)\,\xi(u)\dd u$.
\end{proposition}
\begin{proof}
Take $\E[\cdot]$ in Model~\ref{mod:ch3-hawkes} and use, by the tower property,
$\E[\dd N(u)]=\E[\lambda(u)]\dd u=\xi(u)\dd u$ (Tonelli justifies the interchange).
\end{proof}

\begin{intuitionbox}
Averaging the \emph{stochastic} feedback over all realisations turns it into
\emph{deterministic} feedback: ``on average, today's rate is the baseline plus the
kernel-weighted echo of the average past.'' This equation --- a Volterra equation
of the second kind (Appendix~\ref{app:volterra}) --- is the protagonist of the
entire book.
\end{intuitionbox}

\begin{examplebox}[: total progeny and slow Monte Carlo]
A subcritical cluster's size is Borel-distributed with mean $1/(1-\kappa)$ and
variance $\kappa/(1-\kappa)^3$. Near criticality ($\kappa\to1$) the variance
explodes: a few giant cascades dominate, and Monte-Carlo estimates of means
converge painfully slowly --- a recurring theme when we validate near-critical fits.
\end{examplebox}

\begin{recap}
The Hawkes process is a forest of branching family trees; $\kappa=\int\phi$ is the
mean offspring per event and $\kappa<1$ is stability. It is not Poisson (random
compensator), but its mean intensity solves the Volterra equation
$\xi=s+\phi*\xi$. Solving and exploiting that equation is Part II.
\end{recap}

\exercise{Verify by simulation that the mean offspring per immigrant is
$\kappa/(1-\kappa)$ for $\kappa=0.6$. (The package test
\code{test\_total\_progeny\_equals\_branching\_sum} does this.)}
\exercise{Explain why the cluster representation makes the \emph{mean} intensity
deterministic even though each realisation's intensity is random.}
